% Generated draft: requires mathematical and source review.
\documentclass[11pt]{article}
\usepackage[margin=25mm]{geometry}
\usepackage{fontspec}
\setmainfont{Latin Modern Roman}
\usepackage{amsmath,amssymb,mathtools,unicode-math}
\setmathfont{Latin Modern Math}
\usepackage{enumitem,needspace,xurl,hyperref}
\setlength{\parindent}{0pt}
\setlength{\parskip}{5pt}
\setlength{\emergencystretch}{3em}
\newcommand{\N}{\mathbb{N}}
\newcommand{\Z}{\mathbb{Z}}
\newcommand{\Q}{\mathbb{Q}}
\newcommand{\R}{\mathbb{R}}
\newcommand{\C}{\mathbb{C}}
\newcommand{\F}{\mathbb{F}}
\newcommand{\A}{\mathbb{A}}
\newcommand{\PP}{\mathbb{P}}
\newcommand{\E}{\mathbb{E}}
\newcommand{\eps}{\varepsilon}
\newcommand{\dd}{\,\mathrm{d}}
\begin{document}
\section*{UnsolvedMath ID: 11000044}
\textbf{Source problem code:} catalogue entry
\subsection{Definitions and mathematical statement}
For a surface \(S\) with the source's complexity \(d(S)\geq1\), let \(M\) be the moduli-type space denoted by \(M\) in the preceding discussion, equipped with its Teichm\"uller metric. The conjecture asserts that \(M\) admits no finite-volume Riemannian metric of ``(uniformly bounded) positive scalar curvature'' whose metric space lies in the Teichm\"uller quasi-isometry class. Neither \(d(S)\) nor the exact \(M\) (moduli space, orbifold or prescribed manifold cover) is defined in the excerpt.

A metric \(h\) has finite volume when \(\int_M\dd\operatorname{vol}_h<\infty\). Quasi-isometry means there is a coarsely surjective map \(F\) and constants \(A\geq1,B\geq0\) with
\[A^{-1}d_{\mathrm{Teich}}(x,y)-B\leq d_h(Fx,Fy)\leq A d_{\mathrm{Teich}}(x,y)+B.\]
The source does not say that \(F\) must be the identity. It also omits completeness. The curvature wording needs clarification: pointwise \(\operatorname{Scal}_h>0\), a uniform positive lower bound \(\operatorname{Scal}_h\geq c>0\), an upper bound, and two-sided positive bounds are different conditions. The exact quoted parenthesis is preserved instead of silently choosing one.

The source reports an analogue for \(\Gamma\backslash G/K\) proved by S. Chang [Ch], and proposes that the same method reduce the present conjecture to understanding large-scale Teichm\"uller geometry. For a fixed basepoint \(x_0\in\mathcal M_g\), it defines
\[\operatorname{Cone}(\mathcal M_g)=\lim_{n\to\infty}(\mathcal M_g,n^{-1}d_{\mathrm{Teich}},x_0),\tag{10}\]
in the pointed Gromov--Hausdorff sense, and asserts existence of this limit in its setting. This is an ordinary specified pointed limit, not an ultralimit selected here. The exact geometric hypotheses for the asserted existence and reduction are not supplied, and the conjectural description announced after this setup is missing.

The following background defines the curve complex \(\mathcal C_g\), for \(g\geq2\): its vertices are nontrivial nonperipheral isotopy classes of simple closed curves on \(\Sigma_g\), and a set of \(k+1\) distinct vertices spans a \(k\)-simplex when they admit mutually disjoint representatives. The mapping class group acts simplicially. The source attributes the large-scale discussion to conversations with H. Masur. No unprovided identification of the cone with a quotient or cone of this complex is asserted.
\subsection{Short English statement}
The conjecture forbids certain finite-volume positive-scalar-curvature metrics in the Teichm\"uller quasi-isometry class. Its exact curvature and space conventions are missing; the following tangent-cone and curve-complex setup is preserved without inventing the announced next conjecture.
\subsection{Sources}
\begin{itemize}
\item Local UnsolvedMath registry, permanent ID \(11000044\), exact source code \texttt{catalogue entry}.
\item Farb, \emph{Problems on Mapping Class Groups and Related Topics}, Conjecture 4.6, PDF page 39, equation (10): \url{https://www.math.uchicago.edu/~farb/papers/mcgbook.pdf}. Registry access 2026-07-29; pdf-font-structure extraction.
\item S. Chang [Ch] is cited for the \(\Gamma\backslash G/K\) analogue; the full entry and precise theorem are absent. H. Masur is credited for discussions underlying the large-scale question.
\end{itemize}
\end{document}
